% =====================================================================
%  DANH MỤC TOÁN 11 — tên bộ tài liệu, chương, chuyên đề.
%  Dùng trong file .tex:  \loai{1}\chuong{3}   |   \loai{2}\de{01}   |   \loai{3}\chuyende{1}
% =====================================================================
\dmuc{loai1}{CÁC CHUYÊN ĐỀ CƠ BẢN TOÁN 11}
\dmuc{loai2}{ĐỀ KIỂM TRA TOÁN 11}
\dmuc{loai3}{CHUYÊN ĐỀ ÔN THI TOÁN 11}

% Chương SGK Toán 11 Kết nối tri thức (kiểm lại theo sách đang dạy)
\dmuc{chuong1}{Chương I. Hàm số lượng giác và phương trình lượng giác}
\dmuc{chuong2}{Chương II. Dãy số. Cấp số cộng và cấp số nhân}
\dmuc{chuong3}{Chương III. Các số đặc trưng đo xu thế trung tâm của mẫu số liệu ghép nhóm}
\dmuc{chuong4}{Chương IV. Quan hệ song song trong không gian}
\dmuc{chuong5}{Chương V. Giới hạn. Hàm số liên tục}
\dmuc{chuong6}{Chương VI. Hàm số mũ và hàm số lôgarit}
\dmuc{chuong7}{Chương VII. Quan hệ vuông góc trong không gian}
\dmuc{chuong8}{Chương VIII. Các quy tắc tính xác suất}
\dmuc{chuong9}{Chương IX. Đạo hàm}
\dmucngan{chuong1}{Chương I. Hàm số lượng giác, phương trình lượng giác}
\dmucngan{chuong3}{Chương III. Số đặc trưng của mẫu số liệu ghép nhóm}

% Chuyên đề (đặt tạm, sửa theo ý)
\dmuc{chuyende1}{Chuyên đề 1. Lượng giác}
